% ===== §5 Epistemology II: value rewrites the subject =====
\section{The Political Economy of the Reconfigured Subject}
\label{sec:political}

\noindent
The reflexive structure is now realised under the conditions of a second discipline, in which it assumes a distinct form. In \xref{sec:grammar} it is realised as a formal epistemology of participatory measurement; here it is realised as a political economy. The two sections are deliberately symmetric and deliberately independent: each takes the single abstract structure, the reconfiguration of the dynamics by the phenomenon, of the producer by the product, and develops what its discipline establishes concerning that structure, to the discipline's limit, without reduction of either to the other. Political economy contributes what the formal articulation does not: that the producer reconfigured by the product is a subject, and that the reconfiguration possesses a determinate direction that the geometry of \pinkref{Paper~IX} specifies.

\subsection{The reflexive structure in the register of production}
\label{sec:reflexive-production}

The abstract structure is stated in political-economic terms. A generative system produces; its product reacts upon and reconfigures the system's dynamics. Under the reading of ``produces'' as social production and of ``the system's dynamics'' as the producing subject, the structure yields a thesis that is Marx's: the subject creates value, and the value, in circulation, reacts upon and reconfigures the subject who created it. The product is not inert. Value, once created and set in circulation, does not remain the completed output of a subject who persists unaltered; it returns upon its creator and reconstitutes the creator. This is the reflexive structure of \xref{sec:reflexive} in the social register, and it is the same structure: the appearance reconfigures the movement, the product reconfigures the producer.

Political economy contributes what the formal section does not. The reconfiguration is not neutral; it possesses a direction, and the direction is the distinction the prior paper drew between the vicious and the good cycle. The reconfiguration of the subject by its product is, in Marx, divided, at once the developed form of alienation and the form of self-realisation, and the determination of which obtains is given by the sign of the phase.

\subsection{The negative phase: the alienated reconfiguration}
\label{sec:alienated}

In the negative case the reconfiguration depletes. The worker produces value; the value congeals as capital; and capital, the ``dead labour'' of \emph{Capital}, reacts upon living labour to dominate, deplete, and diminish it \citep{marx1976,marx1959}. The product reconfigures the producer negatively: the subject is subordinated to its own product, diminished by its own creation, and committed to sustaining a circulation that expands as the subject contracts. This is the reflexive structure in the orientation contrary to the good: the reconfiguration obtains, it obtains necessarily, the structure guaranteeing it, but the accumulated phase is extracted and possessed elsewhere, and the creating subject is depleted, a means to a cycle whose increment accrues to another. Capital is the producer reverse-colonised by its product: a reflexive structure in which the reconfiguration of subject by value proceeds at negative phase, the subject depleted in the measure that the circulation is augmented. This is the vicious cycle of the prior paper, the structure of net extraction, of negative holonomy, located at its sharpest point, the production of the subject.

\begin{claim}[Alienation as the negative-phase reconfiguration of the subject]
Alienation is not, fundamentally, the loss of a product, the worker's forfeiture of the produced thing. It is the structural foreclosure of the subject's positive reconfiguration by its own product: the subject creates value, the value reconfigures the subject, but the reconfiguration is a depletion, the phase extracted and possessed elsewhere, such that the creator is reconstituted as diminished rather than augmented by what it creates. The fundamental sense of alienation is the foreclosure of the subject's becoming, through its own creating, the subject it might have become. The fuel-subject of the prior paper, the subject that creates value and is never reached by its return, is, in this ontology, the subject for which the positive reconfiguration is foreclosed: present in the cycle solely as that consumed to drive it.
\end{claim}

\subsection{The positive phase: the generative reconfiguration}
\label{sec:generative-rewrite}

The same reflexive structure admits the contrary orientation, of which Marx also gives an account, most explicitly in the \emph{Grundrisse}: labour as the objectification of human powers, and wealth in its non-alienated form as the developed richness of human needs, capacities, and relations \citep{marx1973}. The product reconfigures the producer positively: the subject creates value, and the circulation of the value returns to augment and develop the subject that created it, the subject becoming, through its own creating, a more developed subject. This is the good cycle of the prior paper, the positive holonomy, located at the production of the subject: a structure in which the reconfiguration of subject by value proceeds at positive phase, the creator augmented by what it creates, the return accruing to the one whose creation it was.

The structure thus yields a thesis stronger than either pole considered alone.

\begin{claim}[The reconfiguration is necessary; its direction is the variable]
That the circulation of value reconfigures its creating subject is not contingent but a structural necessity of the reflexive system, as a quantum measurement necessarily reacts upon its state. The variable is not whether the subject is reconfigured but the direction, the sign of the phase: whether the subject is depleted by its own product (alienation, negative phase, the vicious cycle) or developed by it (self-realisation, positive phase, the good cycle). Settling, possessive, extractive circulation effects the former; non-possessive circulation effects the latter. The apparatus of the prior papers, the good and vicious cycle, the geometric phase, possession and the dark virtue, settling failure, the fuel criterion, is, in this section, reactivated and unified at the central object of political economy: the production of the subject.
\end{claim}

\subsection{The non-circular criterion of the phase sign}
\label{sec:surplus}

The identification of alienation with negative phase and of self-realisation with positive phase requires an independent criterion for the sign, on pain of circularity: were the sign defined by the alienation it is invoked to explain, the identification would be vacuous. The criterion is supplied by a reproduction surplus, defined independently of the phase. Let a cycle of the relational or productive system be considered over one period. Let $G$ denote the gross regeneration effected in the period, the augmentation of the subject's generative conditions: its capacities, relations, and powers of further production. Let $C$ denote the consumption the period exacts of the subject, the depletion of those same conditions in driving the cycle. Define the reproduction surplus as
\[
\sigma \;=\; G - C .
\]
A cycle for which $\sigma > 0$ regenerates the subject's generative conditions in excess of what it consumes: the subject emerges with augmented capacity for further generation. A cycle for which $\sigma < 0$ consumes those conditions in excess of what it regenerates: the subject emerges depleted. A cycle for which $\sigma = 0$ reproduces the subject's conditions without augmentation, the structure of simple repetition. The thesis is then that the sign of $\sigma$ tracks the sign of the holonomy: positive reproduction surplus corresponds to positive phase, negative surplus to negative phase, and zero surplus to the zero-holonomy cycle of mere repetition. The correspondence is asserted as a thesis to be established, not as a definition, and the independent definability of $\sigma$ in terms of regeneration and consumption, without reference to the phase, is what secures the identification of \xref{sec:alienated} and \xref{sec:generative-rewrite} against circularity.

\subsection{The subject as generated within the cycle}
\label{sec:subject-in-cycle}

A further consequence completes the symmetry with \xref{sec:grammar}. The reconfigured subject is not a complete subject that entered production antecedently constituted and was subsequently modified. The subject is generated within the cycle of which it is also the producer. There is no fixed, antecedently given creator of value standing external to the circulation, awaiting the production and the subsequent modification by what it produced; the subject is itself produced, recurrently, in and through the cycle, as the social-reproduction tradition establishes, the subject being reproduced in each period through the labour of care and maintenance on which the prior papers have insisted \citep{federici2004,fraser2017,bhattacharya2017}. The subject is the fixed point of the reflexive structure, not its external premise: that which is produced in the producing, constituted in the constituting. This is the political-economic form of the first consequence of \xref{sec:reflexive}, and of the mother paper's relational ontology, that the subject is generated in and through relation rather than entering it complete. Marx's thesis that in transforming the objective world the human transforms itself, Foucauldian subjectivation, and the recurrent reproduction of the labouring subject are, in their distinct registers, formulations of the single thesis of this ontology: the producer is among the products.

\subsection{The identity of the reconfigured subject as parallel transport}
\label{sec:identity}

A difficulty requires address. If the subject is recurrently reconfigured by the cycle, the sense in which value is ``returned to its creator'' returned to the same subject is in question. The creator augmented by the return is not, subsequent to the reconfiguration, identical to the creator that first created the value, the reconfiguration having reconstituted it. The question is whether ``return to the creator'' thereby loses its sense, there being, strictly, no single creator to which value returns.

The sense is not lost but is specified. This is the central result of the prior paper, the return to the root that is not a return to the origin, the spiral as distinct from the circle, obtaining at the level of the subject. The value returns to its creator, but the creator has, in the returning, been augmented; the return is therefore not to the original subject but to the spirally developed one. ``Creator and beneficiary are identical'' is not the persistence of a self-identical subject but the spiral development of a subject. The geometric formulation specifies the result:

\begin{claim}[Personal identity as parallel transport in the reflexive value-cycle]
The continuity of the subject through the reflexive cycle is not the strict self-identity of an unchanged entity but the parallel transport of a subject around a loop in relational space, a sameness that carries, on return to its starting configuration, a phase difference. Personal identity, on this construal, is parallel transport in the reflexive value-cycle; alienation is the negative-phase displacement of that identity, the non-recognition of itself by the depleted subject; and self-realisation is its positive-phase development, the becoming more itself of the subject through its own creating. The creator to which value returns is the creator parallel-transported around the cycle: the same, and developed, returned to the root and not to the origin. The identity of creator and beneficiary holds in the holonomic sense, the sense of a sameness that has accumulated a phase. The construal is advanced as a structural model of the subject's continuity, and the adequacy of parallel transport as its formal carrier is subject to the cautions of \xref{app:phase}.
\end{claim}

\noindent
This is the limit of the section, the point at which political economy, developed to its boundary, returns the question to the geometry without being absorbed by it. It has established what its discipline establishes: that the reflexive reconfiguration of the producer by the product is a social fact possessing a moral direction; that alienation and self-realisation are the two phases of one cycle, distinguished by the sign of an independently definable reproduction surplus; that the subject is produced within the cycle it produces; and that the return to the creator is a spiral and not a circle, a sameness bearing a phase. It has established this in its own terms, to its own boundary, without dissolution into the formal grammar of \xref{sec:grammar} or legislation by the geometry it employs. Two realisations of one structure, the formal-epistemic and the political-economic, now stand in juxtaposition, recognisably identical in structure and irreducibly distinct in register. The status of the situation in which one abstract structure is realised, under distinct disciplines, in such distinct and mutually irreducible forms, and the cognition of the unbroken structure they share, in the absence of the external standpoint excluded in \xref{sec:onto-refuses}, are the matter of \xref{sec:breaking}.
